\section{The original sphere comparison and its complete exercises}\label{courseIII:integral-homology-and-sphere-recognition:section}

The following complete arguments are retained as re-derivations of their classical prerequisites. Source-specific calculations and their hypotheses are stated explicitly.

\subsection{The actual smooth recognition}
\label{courseIII:integral-homology-and-sphere-recognition:smooth-recognition}

\label{courseIII:integral-homology-and-sphere-recognition:section-10}

The homotopy equivalence of Section 9 identifies a smooth homotopy six-sphere. The smooth statement uses two classical results: the simply connected smooth \(h\)-cobordism theorem and the vanishing of the group \(\Theta_6\) of oriented smooth homotopy six-spheres. The \cref{courseIII:sphere-twists-and-the-original-smooth-six-sphere-comparison:twist-original-smooth-map} now constructs the new interior gluing with its full frame defect, cancels every actual handle and gives the orientation-preserving diffeomorphism to the original six-sphere, homotopic to its original degree-one map. We first construct the precise cobordism to which the smooth theorem applies.

An \(h\)-cobordism from a closed manifold \(M_0\) to \(M_1\) is a compact smooth manifold \(C\) with \(\partial C=M_0\amalg M_1\) for which both inclusions \(M_i\hookrightarrow C\) are homotopy equivalences. It comes with its actual boundary inclusions.

\begin{proposition}
Choose two disjoint orientation-preserving coordinate closed discs \(D_0,D_1\subset X\). Then


```{=latex}
\begin{equation}
C=X\setminus(\operatorname{int}D_0\cup\operatorname{int}D_1)
\label{courseIII:integral-homology-and-sphere-recognition:eq:10.1}
\end{equation}
```


is a simply connected smooth six-dimensional \(h\)-cobordism between its two boundary five-spheres.
\end{proposition}
\begin{proof}
Coordinate discs have collars, so \eqref{courseIII:integral-homology-and-sphere-recognition:eq:10.1} is a compact smooth manifold with the indicated boundary. It is connected: any path between points outside the discs can be perturbed to meet the boundary spheres transversely in finitely many points, and each portion passing through a disc can be replaced by a path in its connected boundary sphere. Van Kampen, applied when adding the two discs back along collars of \(S^5\), identifies \(\pi_1(C)\) with \(\pi_1(X)\). The discs and the collar spheres are simply connected. Thus \(\pi_1(C)=0\).

Excision gives


```{=latex}
\begin{equation}
H_k(X,C;\mathbb Z)=
H_k(D_0,\partial D_0;\mathbb Z)\oplus
H_k(D_1,\partial D_1;\mathbb Z)
=
\begin{cases}\mathbb Z^2,&k=6,\\0,&k\ne6.\end{cases}
\label{courseIII:integral-homology-and-sphere-recognition:eq:10.2}
\end{equation}
```


Here each summand is computed by the relative cell of the disc; its positive generator is its local orientation class. The fundamental class of \(X\) maps to \((1,1)\), since both chosen coordinate discs preserve the orientation. The long exact sequence of the pair and Theorem 7.1 consequently give


```{=latex}
\begin{equation}
H_6(C)=0,\qquad
H_5(C)=\mathbb Z^2/\mathbb Z(1,1)\simeq\mathbb Z,
\qquad H_k(C)=0\quad(1\leq k\leq4).
\label{courseIII:integral-homology-and-sphere-recognition:eq:10.3}
\end{equation}
```


The relative boundary homomorphism sends the first and second standard generators in \eqref{courseIII:integral-homology-and-sphere-recognition:eq:10.2} to their actual boundary sphere classes in \(C\). In the quotient in \eqref{courseIII:integral-homology-and-sphere-recognition:eq:10.3} these are opposite generators. Thus each inclusion \(\partial D_i\hookrightarrow C\) induces an isomorphism on \(H_5\), and also on \(H_0\); all other homology groups on both sides vanish.

Both spaces are simply connected. The included \cref{mapsI:cw:cw-manifolds} proves their CW homotopy type: its formula (2.6) treats the actual punctured cobordism, keeping the original disk radii. The included \cref{mapsI:pf:integral-comparison}, Theorem 5.1, turns these homology isomorphisms into homotopy equivalences. Surjectivity of the degree-\(n+1\) edge first kills the only possible incoming transgression; injectivity in degree \(n\) then excludes a first nonzero fibre group. Its equation (5.12) constructs both original inverse maps through the specified CW comparison of \(C\), retaining the opposite signs of the two boundary generators. Hence both actual boundary inclusions are homotopy equivalences, proving the assertion.
\end{proof}

Section 6 of the \cref{courseIII:realizing-framed-handles-and-middle-cancellation:middle-six-seven-reduction} proves the required smooth product theorem in dimensions six and seven. Its exact map (6.8) identifies this \(C\) with \(\partial D_0\times[0,a_+-a_-]\), relative to its incoming boundary. Here \(a_-,a_+\) are the retained endpoint values of the resulting function. The original coordinate parametrization \(j_0:S^5\to\partial D_0\) and the map \((z,t)\mapsto(j_0(z),(a_+-a_-)t)\) compare this product with \(S^5\times[0,1]\); the inverse divides the second coordinate by the same positive factor and applies \(j_0^{-1}\). The full derivative retains \(Dj_0\) and the factor \(a_+-a_-\). Adding the two original coordinate discs then describes \(X\) as

\begin{equation}
D^6\cup_\phi D^6
\label{courseIII:integral-homology-and-sphere-recognition:eq:10.4}
\end{equation}

for a particular smooth gluing diffeomorphism \(\phi:S^5\to S^5\). The smooth map on the remaining boundary is part of the data in \eqref{courseIII:integral-homology-and-sphere-recognition:eq:10.4}; a product description of the complement does not itself extend \(\phi\) across a disc.

The classical group \(\Theta_6\) uses oriented homotopy spheres modulo oriented \(h\)-cobordism, with connected sum as its operation. The \cref{courseIII:sphere-twists-and-the-original-smooth-six-sphere-comparison:twist-original-smooth-map} gives the complete receiving proof that every oriented smooth homotopy six-sphere is orientation-preservingly diffeomorphic to the standard sphere. Its proof constructs a suitable seven-dimensional product with the specified original boundary, and thus proves the group equality as well. The full application to the original map of Section 9 is given below.

The \cref{courseIII:relative-morse-functions-and-original-handles:relative-morse-attachment} now constructs the actual starting decomposition of this original cobordism. It fixes the original incoming and outgoing collars, retains every Hessian coefficient in its exact coordinate map, and inserts each framed handle with both original parameter radii. Its Sections 4.2--4.4 specify the transverse seam, outward corner smoothing and complete variable-time collar map, including both extreme indices and the final outgoing boundary identification. This proves the finite handle decomposition used below. The \cref{courseIII:rearranging-the-original-framed-handles:arrangement-interchange} is now proved as well. It compares the exact rounded boundary with the original surgery collars, moves complete attaching tubes through the original complement maps, interchanges the actual framed handles and retains the final boundary identification. The \cref{courseIII:morse-trajectories-and-critical-value-lowering:trajectory-ordered-morse} supply the compatible function and the full supported critical-value change, including all original Hessian coefficients and the actual boundary maps. The \cref{courseIII:removing-the-original-zero-and-top-handles:zero-supported-cancellation} is now proved as well, retaining both original boundary collars and the entire presentation map. The \cref{courseIII:creating-an-original-cancelling-handle-pair:pair-original-band} supplies the exact local insertion needed for stabilization, with a single transverse connecting trajectory and the full original-band comparison. Its lower attaching sphere bounds the explicitly framed disk proved in Section 6 of that companion. The \cref{courseIII:framed-handle-slides-and-index-one-removal:slide-index-one-cancellation} now remove index one and its reversed counterpart by the actual supported belt crossing, full loop isotopy, exact common-exterior flow comparison and verified one-intersection cancellation. Section 8.1 of that companion proves the stronger surjectivity hypothesis for the boundary fundamental-group map. Only the original handles of indices two through dimension minus two remain. The \cref{courseIII:integral-handle-chains-and-signed-incidences:handle-six-seven} now proves the original filtration-to-homology map, every incidence sign, the full integer basis comparisons and the exact contractions in dimensions six and seven. The \cref{courseIII:higher-index-framed-slides-and-chain-maps:higher-slide-matrices} now realizes every elementary integral change with its full framed cobordism map and proves both neighbouring differentials using the singular-prism correction. The \cref{courseIII:whitney-moves-in-the-receiving-dimensions:receiving-whitney-flow} proves the full disk, adapted frame and supported move for the three required sheet-dimension pairs, with all original metric, orientation and framing maps retained. The \cref{courseIII:realizing-framed-handles-and-middle-cancellation:middle-six-seven-reduction} constructs the compatible function, proves the exact intersection-to-trajectory comparison and every escape and support hypothesis, and completes the simply connected smooth cobordism reductions in dimensions six and seven.

The \cref{courseIII:morse-cancellation-with-controlled-support:section} now supplies the geometric cancellation step, including the supported critical-value change and the exact product map relative to the lower boundary. Its Section 5 retains the disk radius and all derivatives of the cutoffs, and its second exercise computes the extra critical points produced by an invalid cutoff. This is the receiving cancellation argument. Sections 2--6 of the \cref{courseIII:realizing-framed-handles-and-middle-cancellation:middle-six-seven-reduction} now supply the actual function, unique trajectories, compact supports and finite cancellation order needed to apply it to every middle pair.

The \cref{courseIII:belt-sphere-complements-and-the-whitney-disk:whitney-boundary} now proves the complete framed disk construction for a specified simply connected index-two handle stage. Its Sections 1--3 retain the original complement diffeomorphism, both handle radii and both tangent maps. Section 5 constructs embedded fillings while preserving the entire original collar and avoiding the boundary and corner strata of the previously built collar. Sections 6--7 retain every nonorthogonal frame coefficient, the original metric and both intersection signs. They construct the actual two-corner disk and its full adapted frame. The \cref{courseIII:whitney-move-with-controlled-support:move-receiving-attachments} now proves its receiving ambient operation. It constructs the exact corner-domain and sheet coordinates, preserves all original frame factors, removes the two intersections by a full cutoff flow, and gives the resulting diffeomorphism of attached cobordisms relative to the incoming boundary. The \cref{courseIII:whitney-moves-in-the-receiving-dimensions:receiving-whitney-flow} proves the full disk, adapted frame and supported move for the three required sheet-dimension pairs, with all original metric, orientation and framing maps retained. The \cref{courseIII:realizing-framed-handles-and-middle-cancellation:middle-six-seven-reduction} constructs the compatible function, proves the exact intersection-to-trajectory comparison and every escape and support hypothesis, and completes the simply connected smooth cobordism reductions in dimensions six and seven. The resulting product map retains the original incoming inclusion and its actual outgoing flow map.

The \cref{courseIII:spin-six-and-the-framing-comparison:section} also proves a precise tangent comparison for the actual degree-one homotopy equivalence \(g:S^6\to X\) constructed in Section 9. There is an oriented bundle isomorphism

\begin{equation}
B:g^*TX\longrightarrow TS^6.
\label{courseIII:integral-homology-and-sphere-recognition:eq:10.5}
\end{equation}

It is obtained by lifting the original tangent clutching map to \(SU(4)\), using its exterior-square real representation, and proving that the Euler number divided by two classifies oriented rank-six bundles on \(S^6\). Both bundles in \eqref{courseIII:integral-homology-and-sphere-recognition:eq:10.5} have Euler number two. The proof retains the determinant contribution before using its specified trivialization. The actual stabilization map and its inverse are

\begin{equation}
\begin{aligned}
(p;v,a)&\longmapsto(p;B_pv+ap),\\
(p;w)&\longmapsto
\bigl(p;B_p^{-1}(w-\langle w,p\rangle p),\langle w,p\rangle\bigr).
\end{aligned}
\label{courseIII:integral-homology-and-sphere-recognition:eq:10.6}
\end{equation}

Thus \(g^*TX\oplus\mathbf1_{\mathbb R}\) is trivial, and pulling back by the homotopy inverse gives a stable framing of the original \(TX\oplus\mathbf1_{\mathbb R}\). Its rank-six Euler class has not disappeared. A stable framing alone does not compute a framed bordism class or prove \(\Theta_6=0\).

The \cref{courseIII:the-degree-six-milnor-cobar-calculation:six-cobar-survival} now proves the complete degree-six cohomology calculation for the original Milnor Hopf algebra at two and its vanishing at every odd prime. Its proof retains every word length, both full correcting cochains, the source filtration comparison and every odd-prime sign. The actual coefficient comparisons and convergence to the specified arithmetic completions are proved in the subsequent linked companions. The \cref{courseIII:the-original-integral-stable-group-and-its-completion-maps:integral-degree-six-answer} now proves finite generation of the original sphere groups, the exact suspension range and the full unit comparison with all prime completions. It proves \(\pi_6(\mathbb S)=\mathbb Z/2\mathbb Z\) and that the original 2-completion unit is an isomorphism in degree six. The \cref{courseIII:the-original-six-dimensional-framing-invariant:framed-six-original-nullbordism} constructs the full quadratic invariant, proves framed bordism invariance and gives value zero for the original framed homotopy six-sphere, with an actual framed nullbordism. The \cref{courseIII:sphere-twists-and-the-original-smooth-six-sphere-comparison:twist-original-smooth-map} now constructs the new interior gluing with its full frame defect, cancels every actual handle and gives the orientation-preserving diffeomorphism to the original six-sphere, homotopic to its original degree-one map. The algebraic theorem alone does not establish the smooth sphere-group calculation.

The \cref{courseIII:framed-manifolds-and-the-pontryagin-thom-map:pt-bijection} now proves both geometric constructions and their exact inverse comparison. It constructs the original normal frame from the specified tangent framing, retains the full tube and Gram matrix, recovers the original frame at a regular inverse image, and proves the relative homotopy-to-bordism map. Its \cref{courseIII:framed-manifolds-and-the-pontryagin-thom-map:pt-stabilization-and-original-framing} shows exactly how a proved zero stable class would yield a bounding manifold with the given boundary framing. The homotopy of the full comparison involution uses the original complex tangent operator. This completes the geometric correspondence. The \cref{courseIII:the-original-six-dimensional-framing-invariant:framed-six-original-nullbordism} now calculates the original stable class as zero and obtains its framed nullbordism.

The \cref{courseIII:coefficient-rings-and-the-original-degree-six-class:coefficient-original-h28} now proves the full source coefficient and grading comparisons. It calculates the entire original polynomial-coefficient cobar group in bidegree \((2,8)\), proves that this group has no coefficient torsion, and gives the exact map of its generator to the classical degree-six class. All other lengths retain an explicit bound on their coefficient torsion. These algebraic results feed the actual cooperation and integral comparisons proved in the linked companions below. The \cref{courseIII:the-original-six-dimensional-framing-invariant:framed-six-original-nullbordism} now completes the geometric value calculation and supplies the framed nullbordism.

The \cref{courseIII:the-adams-tower-and-exact-inverse-limits:adams-actual-convergence} constructs the ordinary-cohomology spectrum from its full finite cocycles and cochain spaces, proves every stage group finite with its exact prime-power bound, and proves convergence to the specified homotopy limit with a complete separated filtration. Its \cref{courseIII:the-adams-tower-and-exact-inverse-limits:adams-degree-five-cobar} proves all-length vanishing needed by the outgoing degree-six differentials. The \cref{courseIII:the-actual-mod-two-cooperation-comparison-through-degree-eight:actual-eight-bijectivity} identifies the actual cooperation algebra through degree eight. The original geometric framing chapter proves the value zero and its framed nullbordism.

The \cref{courseIII:the-original-adams-tower-and-arithmetic-completion:arithmetic-adams-equivalence} now constructs the unit-preserving equivalence from the original Moore-tower limit to this Adams limit. The proof covers every bounded-below spectrum and gives the exact finite transition bound \(2m(n-b+1)\). Its \cref{courseIII:the-original-adams-tower-and-arithmetic-completion:arithmetic-homotopy-groups} retains both derived terms, the original unit sign, the full torsion transitions and the complete degree-zero extension. In particular \(\pi_0C_s=\mathbb Z/p^{s+1}\), so the stage exponent is sharp. The original geometric framing chapter proves that value zero, with its complete geometric invariant.

The \cref{courseIII:the-actual-mod-two-adams-resolution-and-cooperation-maps:cooperation-chain-comparison} now constructs both cross-product coordinate maps, the actual cofree Adams resolution and both chain maps and inverse homotopies comparing it with the actual cooperation cobar complex. Its \cref{courseIII:the-actual-mod-two-adams-resolution-and-cooperation-maps:cooperation-original-diagonal} identifies the generator tied to the original integer two, proves every degree-two reduced coproduct zero, and shows that the degree-six calculation at every word length uses cooperation degrees through eight. The \cref{courseIII:the-actual-mod-two-cooperation-comparison-through-degree-eight:actual-eight-bijectivity} supplies the required polynomial comparison through degree eight; the \cref{courseIII:the-original-six-dimensional-framing-invariant:framed-six-original-nullbordism} now supplies that value and its actual framed nullbordism.

The \cref{courseIII:the-exact-characteristic-two-coordinate-calculation:formal-two-linearization} now supplies the all-ring formal-series calculation and both inverses. The \cref{courseIII:the-original-projective-coordinate-and-actual-cooperations:projective-actual-coproduct} constructs the polynomial-to-cooperation morphism, proves its full coproduct and antipode comparison, and identifies its \cref{courseIII:the-original-projective-coordinate-and-actual-cooperations:projective-original-two}. The \cref{courseIII:the-actual-mod-two-cooperation-comparison-through-degree-eight:actual-eight-bijectivity} now proves bijectivity of this original morphism in every cooperation degree from zero through eight. Its \cref{courseIII:the-actual-mod-two-cooperation-comparison-through-degree-eight:actual-eight-original-suspension} determines the full graded dimensions, and its \cref{courseIII:the-actual-mod-two-cooperation-comparison-through-degree-eight:actual-eight-nonzero-powers} proves every power of the actual first generator nonzero. The \cref{courseIII:the-actual-mod-two-cooperation-comparison-through-degree-eight:actual-eight-receiving-cobar} identifies the entire degree-six line and proves \(\pi_6(\mathbb S^{\wedge}_2)=\mathbb F_2\), with the original permanent-cycle representative. The \cref{courseIII:the-original-odd-prime-cooperations-and-degree-six-adams-calculation:odd-actual-adams} now proves the full signed bar and original suspension calculation, the actual cooperation degrees needed by stem six, and the signed comparison with the original Adams tower. It proves \(\pi_6(\mathbb S^{\wedge}_p)=0\) for every odd prime. The \cref{courseIII:the-original-integral-stable-group-and-its-completion-maps:integral-degree-six-answer} now proves finite generation of the original sphere groups, the exact suspension range and the full unit comparison with all prime completions. It proves \(\pi_6(\mathbb S)=\mathbb Z/2\mathbb Z\) and that the original 2-completion unit is an isomorphism in degree six. The \cref{courseIII:the-original-six-dimensional-framing-invariant:framed-six-original-nullbordism} constructs the full quadratic invariant, proves framed bordism invariance and gives value zero for the original framed homotopy six-sphere, with an actual framed nullbordism. The \cref{courseIII:sphere-twists-and-the-original-smooth-six-sphere-comparison:twist-original-smooth-map} now constructs the new interior gluing with its full frame defect, cancels every actual handle and gives the orientation-preserving diffeomorphism to the original six-sphere, homotopic to its original degree-one map.

\begin{theorem}
There is an orientation-preserving smooth diffeomorphism \(d:S^6\to X\), with \(d\simeq g\) and \(d^{-1}\simeq h_X\) for the original maps of Section 9.
\end{theorem}
\begin{proof}
The original stable tangent frame in \eqref{courseIII:integral-homology-and-sphere-recognition:eq:10.6} enters the full Pontryagin–Thom comparison and geometric framing proof cited above. For the whole six-dimensional framed bordism group, the \cref{courseIII:sphere-twists-and-the-original-smooth-six-sphere-comparison:twist-general-stable-framing} proves the complete tangent-to-normal framing homotopy without assuming a complex structure; this is an independent bundle calculation before the sphere argument. The latter's Sections 8–9 prove framed bordism invariance and identify the original sixth framed bordism group with \(\mathbb F_2\) through its quadratic invariant. The original \(H_3(X;\mathbb F_2)=0\) makes that invariant zero; injectivity therefore yields an actual framed nullbordism \(W\) with its specified boundary identification \(\iota_X:X\to\partial W\).

Remove a coordinate ball \(\chi(D^7_\rho)\) and keep its parametrization \(j(q)=\chi(\rho q)\), including \(Dj_q(v)=D\chi_{\rho q}(\rho v)\). The framed interior surgeries and exact index reduction in that geometric chapter produce a two-connected cobordism with only indices three and four. The sphere-twist companion, Section 4, proves that its actual common level \(P\) has a free unimodular alternating middle group, with two primitive integral Lagrangians: the lower belts \(b_i\) and the upper attaching classes \(a_j\). Both have the same rank \(r\). Its exact original signed matrix is


```{=latex}
\begin{equation}
D^{\rm old}_{ij}=\sigma_i\lambda(b_i,a_j),\qquad
 H_4(C',S^6)=\ker D^{\rm old},\quad
 H_3(C',S^6)=\operatorname{coker}D^{\rm old}.
 \label{courseIII:integral-homology-and-sphere-recognition:eq:10.7}
\end{equation}
```


The signs \(\sigma_i\) are those of the original ordered handle charts; no acyclicity of this old filling is presumed.

Sections 1–3 of the \cref{courseIII:sphere-twists-and-the-original-smooth-six-sphere-comparison:twist-actual-cancellation} construct the complete radius-dependent tube twist, prove its signed integral action and give a terminating integral algorithm. It supplies an actual orientation-preserving diffeomorphism \(\varphi:P\to P\) carrying the upper Lagrangian onto an integral complementary one. This \(\varphi\) acts on the six-dimensional middle level and is distinct from the five-sphere gluing map \(\phi\) in \eqref{courseIII:integral-homology-and-sphere-recognition:eq:10.4}. Reglue the two original seven-dimensional pieces by


```{=latex}
\begin{equation}
C_\varphi=(C_-\sqcup C_+)/
             \bigl(i_+(p)\sim i_-(\varphi(p))\bigr).
 \label{courseIII:integral-homology-and-sphere-recognition:eq:10.8}
\end{equation}
```


Its complete smooth seam is proved in Section 5 there. Both outer boundary maps remain fixed. If \(\tau_\pm\) are the original piece frames expressed in directed collar coordinates and a common stable rank \(7+k\), their exact new comparison is


```{=latex}
\begin{equation}
A_\varphi(p)=\tau_{-,\varphi(p)}
                (D\varphi_p\oplus1\oplus I_k)\tau_{+,p}^{-1}.
 \label{courseIII:integral-homology-and-sphere-recognition:eq:10.9}
\end{equation}
```


Thus the construction retains the original frames without assuming that the new gluing preserves them.

In a dual basis \(\lambda(b_i,f_j)=\delta_{ij}\), put \(F_i=\sigma_i f_i\) and write \(\varphi_*a_j=\sum_iF_iU_{ij}\). The matrix \(U\) is unimodular because both collections are bases of the same primitive subgroup. The original handle-to-intersection comparison gives the new boundary matrix \(D^\varphi=U\). The full framed slides realize \(U^{-1}\), with each orientation sign recorded as a signed comparison of the original geometric core. Supported Whitney moves in the actual dimension pair \((3,3)\) remove the opposite-sign intersections. The level before the lower stage is the simply connected incoming \(S^6\), so the proved belt-complement, disc and rank-four framing hypotheses all hold. Exactly one transverse geometric intersection remains for each matched pair. Section 6 of that companion applies the included actual Morse realization, checks escape and disjoint compact supports, and cancels all pairs.

The resulting function on \(C_\varphi\) has no critical point. With its retained endpoint values \(\ell_-<\ell_+\), let \(T=\ell_+-\ell_-\) and \(Y=V/df_{\rm end}(V)\), where the original positive field \(V\) matches both boundary collars. The full product and outgoing maps are


```{=latex}
\begin{equation}
\begin{split}
 \Psi(q,t)&=\operatorname{Fl}^{Y}_t(j(q)),\qquad 0\le t\le T,\\
 d(q)&=\iota_X^{-1}\Psi(q,T),\\
 d^{-1}(x)&=j^{-1}\operatorname{Fl}^{Y}_{-T}(\iota_X(x)).
 \end{split}
 \label{courseIII:integral-homology-and-sphere-recognition:eq:10.10}
\end{equation}
```


Their exact inverse identities, smoothness at the collars and orientation are proved there from \(df_{\rm end}(Y)=1\). In particular \(d_*[S^6]=[X]\). The original first Hurewicz isomorphism in Section 9 sends both \([d]\) and \([g]\) to \([X]\); its injectivity, with the unique basepoint change supplied by \(\pi_1X=0\), gives \(d\simeq g\). Composing with the retained inverse homotopies gives \(d^{-1}\simeq h_X\).

Finally retain the old tangent isomorphism \(B\) from \eqref{courseIII:integral-homology-and-sphere-recognition:eq:10.5}. Finite projection transport along this homotopy supplies \(P:g^*TX\to d^*TX\). The exact comparison with the new derivative is


```{=latex}
\begin{equation}
E_q=B_qP_q^{-1}Dd_q,\qquad Dd_q=P_qB_q^{-1}E_q .
 \label{courseIII:integral-homology-and-sphere-recognition:eq:10.11}
\end{equation}
```


Every map is invertible; no claim \(E=1\) or nullhomotopy of \(E\) is used. This proves the theorem with both original maps, their full tangent comparison, both filling matrices and the stable-frame defect retained.
\end{proof}

The same receiving proof applies to every oriented smooth homotopy six-sphere, using the proved \(\pi_5SO(7)=0\) to obtain its original stable frame. It therefore gives \(\Theta_6=0\) without an assumption of a complex structure. The sphere-twist companion contains two additional fully solved matrix examples, including old torsion and the exact new gluing.

\textbf{Status of the classical prerequisite.} The complete receiving proofs are included and have received author self-checking and supplementary finite algebra checks. The bounded original-author Benedetti reading supplies historical context and explicitly does not supply the whole \(h\)-cobordism proof. The complete geometric arguments instead appear in the linked handle, Whitney, Morse, framing and sphere-twist companions. No independent mathematical review or exhaustive literature reading is claimed.

\subsection{Four worked exercises}
\label{courseIII:integral-homology-and-sphere-recognition:solved-exercises}

\label{courseIII:integral-homology-and-sphere-recognition:section-11}

\subsubsection{Exercise 11.1. Detecting the exact finite index rather than its rational rank}
\label{courseIII:integral-homology-and-sphere-recognition:exercise-11.1}

In the order-four case use the original invariant basis

\begin{equation}
h=\gamma\eta_2,\qquad q=uw+6\gamma\delta,
\label{courseIII:integral-homology-and-sphere-recognition:eq:11.1}
\end{equation}

whose pairing on the torus is \(\begin{pmatrix}0&2\\2&12\end{pmatrix}\).
Knowing that the specialization image has index two and does not contain \(h\), recover the image lattice and its pairing on \(S_2\).

\textbf{Solution.} An index-two sublattice is the kernel of a nonzero map from \(\mathbb Z h\oplus\mathbb Zq\) to \(\mathbb Z/2\). Since it does not contain \(h\), that map takes \(h\) to \(1\). Its value on \(q\) is \(a\in\{0,1\}\). Thus the sublattice has basis \(2h,ah+q\). Pairing these vectors on the torus and dividing by the degree-four covering gives

\begin{equation}
\frac14
\begin{pmatrix}2&0\\a&1\end{pmatrix}
\begin{pmatrix}0&2\\2&12\end{pmatrix}
\begin{pmatrix}2&a\\0&1\end{pmatrix}
=
\begin{pmatrix}0&1\\1&a+3\end{pmatrix}.
\label{courseIII:integral-homology-and-sphere-recognition:eq:11.2}
\end{equation}

The original affine descent calculation (2.10a) applies to the proposed class \(a h+q\): its lifted action would require the integer coefficient \(\ell_3=-(a+3)/2\). Thus \(a+3\) is even, forcing \(a=1\). The explicit lift in the companion proves existence as well as necessity, so the exact image is

\begin{equation}
\mathbb Z(2\gamma\eta_2)\oplus
\mathbb Z(\gamma\eta_2+q)
=\{a'\gamma\eta_2+b'q:a'\equiv b'\pmod2\}.
\label{courseIII:integral-homology-and-sphere-recognition:eq:11.3}
\end{equation}

Its pairing in the displayed basis is \(\begin{pmatrix}0&1\\1&4\end{pmatrix}\). This keeps the original \(q\) and does not replace it by a changed basis vector. The cokernel is generated by the class of \(\gamma\eta_2\) and is exactly \(\mathbb Z/2\).

\subsubsection{Exercise 11.2. Recovering the normalization groups over the integers}
\label{courseIII:integral-homology-and-sphere-recognition:exercise-11.2}

For the constant-sheaf normalization complex, compute the entire cohomology of its degree-zero global row

\begin{equation}
\mathbb Z\xrightarrow{0}\mathbb Z^3
\xrightarrow{D}\mathbb Z^2,\qquad
D(a,b,c)=(a-b+c,a-b+c).
\label{courseIII:integral-homology-and-sphere-recognition:eq:11.4}
\end{equation}

Do the corresponding calculation for the degree-two row of Section 3.

\textbf{Solution.} The first group survives as \(\mathbb Z\). The kernel of \(D\) has basis \((1,1,0)\), \((-1,0,1)\): the equation gives \(a=b-c\), and these two coordinates realize every solution uniquely. Thus its degree-one cohomology is \(\mathbb Z^2\). The image is the full diagonal subgroup, because \(D(1,0,0)=(1,1)\). The map \((x,y)\mapsto y-x\) identifies its cokernel with \(\mathbb Z\). All these groups are free.

In the degree-two row the map sends
\(L=aH-b_1E_1-b_2E_2-b_3E_3\) to

\begin{equation}
(s,-s,s),\qquad s=b_1+b_2+b_3-a.
\label{courseIII:integral-homology-and-sphere-recognition:eq:11.5}
\end{equation}

The scalar map \(s:\mathbb Z^4\to\mathbb Z\) is onto, since \(s(-H)=1\). Its kernel is freely parametrized by \(b_1,b_2,b_3\), with \(a=b_1+b_2+b_3\). Hence the row kernel is \(\mathbb Z^3\). The vector \((1,-1,1)\) is primitive; the quotient is explicitly \(\mathbb Z^2\) via

\begin{equation}
(x,y,z)\longmapsto(x+y,z+y).
\label{courseIII:integral-homology-and-sphere-recognition:eq:11.6}
\end{equation}

Its kernel is exactly the subgroup generated by \((1,-1,1)\), and it is onto by choosing \(y=0\). The degree-four normalization class contributes a further \(\mathbb Z\). Combining the rows with their distinct total degrees gives the ranks \(1,2,4,2,1\) of \(H^*(W;\mathbb Z)\). The spectral-sequence filtration cannot introduce torsion: each successive quotient is free, and a short exact sequence with free quotient splits by lifting a basis.

\subsubsection{Exercise 11.3. Reversing the base orientation without losing a sign}
\label{courseIII:integral-homology-and-sphere-recognition:exercise-11.3}

Retain the original classes \(\theta=12\gamma,\ a=2q,\ b=2\gamma uw\), but write the base generator as \(\omega'=-\omega\). Determine all three differentials in that basis by the product equations.

\textbf{Solution.} The first differential is \(d_2\theta=-\omega'\). Write
\(d_2a=r'\omega'[\delta]\) and
\(d_2b=s'\omega'[\gamma\delta]\).
The unchanged identity \(\theta a=12b\), with \(|\theta|=1\), gives

\begin{equation}
12s'=-24-12r',\qquad s'=-2-r'.
\label{courseIII:integral-homology-and-sphere-recognition:eq:11.7}
\end{equation}

The unchanged identity \(ab=0\) and the original exterior signs
\(\delta b=-2\operatorname{vol}\), \(a\gamma\delta=2\operatorname{vol}\)
give

\begin{equation}
-2r'+2s'=0.
\label{courseIII:integral-homology-and-sphere-recognition:eq:11.8}
\end{equation}

Thus \(r'=s'=-1\). All three coefficients change sign together. Their absolute values remain one, so all three maps remain integral isomorphisms onto their targets. This calculation also shows why reversing only the first sign while keeping the other two would violate the multiplicative spectral sequence.

\subsubsection{Exercise 11.4. Why the Euler characteristic does not close the calculation}
\label{courseIII:integral-homology-and-sphere-recognition:exercise-11.4}

Produce a closed simply connected oriented six-manifold with Euler characteristic two but with nonzero second and third integral homology. Compute those groups and compare the precise missing information with Theorem 7.1.

\textbf{Solution.} Take

\begin{equation}
M=(S^2\times S^4)\mathbin{\#}(S^3\times S^3).
\label{courseIII:integral-homology-and-sphere-recognition:eq:11.9}
\end{equation}

Each product is simply connected. Removing a coordinate disc does not change its fundamental group: adding it back uses a simply connected \(S^5\) collar, so Van Kampen identifies the two groups. Gluing the punctured products along that same simply connected collar shows \(\pi_1(M)=0\).

For \(1\leq k\leq4\), removal of a disc preserves \(H_k\) by the pair sequence and excision, whose only nonzero relative group is in degree six. The Mayer--Vietoris sequence for the connected sum has overlap \(S^5\), with zero homology in degrees one through four; its degree-zero maps also show connectedness. It follows that

\begin{equation}
H_2(M)=\mathbb Z,\qquad H_3(M)=\mathbb Z^2,\qquad H_4(M)=\mathbb Z.
\label{courseIII:integral-homology-and-sphere-recognition:eq:11.10}
\end{equation}

The product groups here follow from their product cells: \(S^2\times S^4\) has cells in degrees \(0,2,4,6\), and \(S^3\times S^3\) has cells in degrees \(0,3,3,6\), with zero cellular boundaries. Orientation and Poincaré duality give \(H_6(M)=\mathbb Z\), \(H_5(M)=H^1(M)=0\); the remaining groups are as stated. Therefore

\begin{equation}
\chi(M)=1+1-2+1+1=2.
\label{courseIII:integral-homology-and-sphere-recognition:eq:11.11}
\end{equation}

Its nonzero \(H_2\) and \(H_3\) prevent it from being a homology sphere. In the original \(X\), it is the three primitive integral transgressions, the two parabolic vanishing calculations and the universal-coefficient argument that remove those groups and their possible torsion. The matching Euler characteristic alone supplies none of these maps.
